\section{正则化与泛化}

\subsection{过拟合问题}

在训练神经网络时，一个核心挑战是\textbf{过拟合}（Overfitting）：模型在训练数据上表现很好，但在从未见过的测试数据上表现很差。这意味着模型"记住"了训练数据中的噪声和特殊模式，而没有学到真正可泛化的规律。

\begin{importantbox}{欠拟合、适度拟合与过拟合}
\begin{itemize}
\item \textbf{欠拟合}（Underfitting）：模型太简单，无法捕捉数据中的基本模式，在训练集和测试集上都表现差
\item \textbf{适度拟合}（Good Fit）：模型恰当地学习了数据中的核心模式，在训练集和测试集上都表现好
\item \textbf{过拟合}（Overfitting）：模型太复杂，拟合了训练数据中的噪声，训练集表现好但测试集表现差
\end{itemize}
\end{importantbox}

\subsection{正则化技术}

为了缓解过拟合，需要使用\textbf{正则化}（Regularization）技术。Amini 介绍了两种主要方法：

\textbf{1. Dropout}

Dropout 的核心思想极其简单：在训练过程中，以概率 $p$ 随机"关闭"（设为零）一部分神经元。

\begin{itemize}
\item 每个 mini-batch 的前向传播中，随机选取约 $p\%$ 的神经元使其输出为零
\item 这迫使网络不依赖于任何单个神经元，而是学习更加分布式和鲁棒的特征表示
\item 测试时不使用 Dropout，但需要按比例缩放输出
\end{itemize}

\begin{lstlisting}[caption={Dropout 的代码实现}]
# TensorFlow
tf.keras.layers.Dropout(rate=0.5)

# PyTorch
nn.Dropout(p=0.5)
\end{lstlisting}

\textbf{2. Early Stopping}

Early Stopping 监控模型在验证集上的表现：当验证损失不再下降（甚至开始上升）时，停止训练。这利用了一个关键洞察：过拟合通常发生在训练后期，此时模型开始"记忆"训练数据而非学习泛化规律。

\begin{warningbox}{正则化不是可选的}
在实际应用中，不使用正则化的深度神经网络几乎一定会过拟合，特别是当模型参数数量远大于训练样本数量时（这在现代深度学习中很常见）。Dropout 和 Early Stopping 应作为标准实践，而非可选的附加项。
\end{warningbox}

\subsection{本章小结}

过拟合是深度学习训练中的核心挑战。Dropout 通过随机关闭神经元迫使网络学习鲁棒特征，Early Stopping 通过监控验证集损失在恰当时机停止训练。二者结合可以显著提升模型的泛化能力。

